% Korean edition of figs/tests.tex.
\fig[wide]{tests}{실서버 테스트에서 드러난 결함 2개를 고친 뒤 모든 테스트가 통과했다}[목표(\figref{goals})별 래퍼 테스트: 2026년 11월 24일 09:27\textasciitilde 10:34 KST 대회 서버 대상 실서버 테스트(포트폴리오 하나)와 오프라인 테스트]{%
{\tabfont\begin{tabularx}{\linewidth}{>{\RaggedRight\arraybackslash}p{16mm} >{\RaggedRight\arraybackslash}p{25mm} >{\RaggedRight\arraybackslash}p{47mm} >{\RaggedRight\arraybackslash}X}
\thd{목표} & \thd{테스트} & \thd{검사한 것} & \thd{결과}\\
G1 & \code{transport} & 첫 요청, 두 번째 클라이언트의 첫 요청, 재사용 연결의 비용. \code{Date} 헤더로 구한 시계 오차 & 새 프로세스의 첫 요청 1.2\textasciitilde 1.3\,s. 두 번째 클라이언트의 첫 요청 130\ms. 재사용 31\textasciitilde 39\ms. 서버 시계 빼기 PC 시계 $-2.98$\,s로 \ref{sec:auth}절의 3\,s와 일치\\\rowsep
G1 & \code{auth} & 두 번 실행하며 로그인, 쿠키 저장, 갱신 & 첫 실행은 로그인하고 쿠키 묶음을 \code{~/.config/pytimefolio/session.json}에 권한 0600으로 저장. 두 번째 실행은 로그인 없이 토큰 갱신\\\rowsep
G2, G5 & \code{rest} & 포트폴리오 규칙. 연결 하나에서 동시 조회. 종목코드 검사 & 규칙으로 운용 중인 포트폴리오를 선택. \code{Session}, \code{Universe}(2,648), \code{Trades}(64), \code{Orders}(그날 오전 5건)를 HTTP/2 연결 하나에서 동시에 가져옴. \code{A005930}의 \code{ValidateProduct}는 경고 없음\\\rowsep
G2 & \code{hub}, \code{state} & 3\,s마다 폴링한 연결 하나. 델타 병합 & 스냅샷 1개 뒤 델타 20개, 끊김 없음. 병합 상태가 \code{t}가 같은 스냅샷과 15회 중 15회 일치(\ref{sec:feed}절의 테스트)\\\rowsep
G3 & \code{orders}, 매수 & 삼성전자 0.1\%, 지정가 150,700원, 시장가 188,400원보다 20\% 낮음 & \code{AddOrder} 후 1\textasciitilde 3\,s 만에 피드에서 찾음: Id 1650412, \code{Working}, \code{qty2ord} 4. 3\,s 뒤 정지. \code{wait()}은 정지, 4주 중 0주 반환\\\rowsep
G3 & \code{orders}, 매도 & LIG넥스원 0.1\%, 지정가 427,000원, 시장가 355,500원보다 20\% 높음 & 주문 레코드 \code{ls} = L, \code{ex} = X, \code{sgn} = $-1$, \code{qty2ord} 2. 정지. 결과 정지, 2주 중 0주\\\rowsep
\hl G3 & \code{plan()}, 합계 & 삼성전자 50\%, \code{weiLimit} 40\% 대비 & 40\%로 상한 적용. water-filling이 초과분을 다른 보유 종목 22개에 분산. 목표 합계 139\%를 \code{plan()}이 받아들임: 결함 1\\\rowsep
G3 & \code{plan()}, 1주(\ref{sec:oneshare}) & \code{min_wei} 0.01로 작은 주문 3건, 09:52 & SK하이닉스 +0.03\%는 경고와 함께 뺌(1주 0.1413\%). 삼성전자 +0.1\%와 리노공업 +0.06\%는 유지\\\rowsep
\hl G3 & \code{guidelines()}(\ref{sec:guidelines}) & 피드 필드와 \code{Violations} 응답, 09:52 & 서버는 위반 0건인데 IT가 섹터 한도를 넘은 것으로 표시: 결함 2. 고친 뒤 서버 위반 0건. 회전율 4.8\%로 \code{ok} False\\\rowsep
G3, G5 & \code{plan()}, 섹터(\ref{sec:guidelines}) & 목표 3가지, 09:55 & 비중 그대로: 주문 0건. 산업재 +6\%p: 거부. 산업재 +3\%p: 주문 2건(\figref{sectortest})\\\rowsep
G2 & \code{live_fields} & 10:33에 60\,s 동안 델타마다 \code{bs_updated}. 주문 없음 & 델타 20개 중 20개에 \code{bs_updated}가 딕셔너리로 옴. \code{na}가 844,505,167원에서 844,076,987원으로 움직여 순자산이 최신으로 유지됨\\\rowsep
G3 & \code{live_rounding}(\ref{sec:oneshare}) & 삼성전자 0.04\%(1.76주어치), 시장가보다 20\% 낮은 지정가. 즉시 정지 & 두 번 모두 \code{qty2ord} 1(Id 1650693 약 10:26, Id 1650719 10:33:45), 0주 체결: 서버는 내림함\\\rowsep
G3 & 수정된 코드, 실서버 & 리뷰 뒤 바꾼 코드로 주문을 넣고 정지: 크기를 0.01\%로 맞춤. 서버가 받았을 수 있는 주문은 재전송 안 함 & 주문 레코드에 \code{wei}가 보낸 그대로 저장(0.04, 0.1, 0.3 확인). 주문 1650719는 피드에서 \code{Accepted}로 찾았고, 즉시 정지했으며, \code{Done}으로 끝남. 10:34 카운터: 새 주문 9건과 정지 7건, 80건 중 16건 사용\\\rowsep
G1\textasciitilde G3, G5 & 오프라인 & 단위 테스트 31개(순수 9개, 모의 HTTP로 실패 경로 22개) & 모두 통과\\
\end{tabularx}}}{시장가는 주문 시점 피드의 가격이다. 테스트 주문은 체결되지 않도록 시장가에서 20\% 떨어진 지정가였다. 강조 행이 결함 2개다. 결함 1: water-filling은 합계를 유지하므로 139\% 목표가 상한 적용 뒤에도 139\%였다. 이제 \code{plan()}은 목표 합계가 \code{max_total}(기본 100\%)을 넘으면 \code{ValueError}를 던진다. 결함 2: 섹터 한도를 섹터 코드 \code{sec}가 아니라 \code{msecs} 레코드 \code{Id}로 찾아 IT를 다른 섹터의 한도와 비교했다. \code{sec}로 찾게 고쳤다. 09:52의 \ref{sec:oneshare}절 테스트와 \ref{sec:guidelines}절 테스트는 \code{plan()}과 \code{guidelines()}만 호출했고 주문은 보내지 않았다. 오프라인 순수 테스트: 상태 병합(문자열 키를 \code{int}로 변환, 바뀌지 않은 필드 유지, 널 응답, 추가·삭제 레코드, 스냅샷 교체), water-filling, 목표-주문 변환, 09:18에 캡처한 요청과 주문 본문 비교, \code{sector_violations}, KRX 호가 단위. 실패 경로(\code{httpx.MockTransport}와 가짜 객체): 상태 코드 401, 403, 429, 400, 502를 예외로 변환. 전송 계층은 연결 실패 때 GET만 한 번 재시도하고 POST는 재시도하지 않으며, \code{AddOrder}는 연결이 맺어지지 않았을 때만 한 번 더 보냄. 읽기 오류는 보낸 것으로 셈. 동시 \code{token()} 호출 5건이 \code{SilentRefresh} 1건. 401이면 한 번 갱신. 손상된 세션 파일이면 로그인. 크기 맞춤. 진행 중 주문은 다시 주문하지 않음. 하루 한도와 동기화. 연결 실패는 세지 않음. 동시 쓰기 3건이 최소 1\,s 간격으로 서버에 도달(1.002\,s 측정). 403과 429 각각이 모든 쓰기를 멈추며, 서버가 다시 200을 줘도 마찬가지. \code{AddOrder} 응답 유실 시 재전송 안 함. 낡은 피드는 주문과 \code{wait()}을 막음. 피드 신선도. 실서버 테스트 뒤 추가: 재시작해도 쿠키 만료 유지. 재전송이 실패해도 검색하며 세 번째 전송 없음. 같은 주문을 두 번 넣으면 ID 2개. 주문 하나에 맞는 새 레코드 2개면 \code{AmbiguousState}. 읽기에서 403이면 쓰기 중단. 서버 핑은 응답 마감을 늘리지 않음. 차단된 피드는 재연결을 멈춤. 이 테스트들은 CI에서 Python 3.11\textasciitilde 3.13으로 돌며, \code{tests/live}는 돌지 않는다.}{DART Quant Team; pytimefolio 0.1.0의 \code{tests/}와 \code{tests/live/}}
